% TODO make hw stylesheet
\documentclass[11pt]{scrartcl}
\usepackage[a4paper, total={6in, 8in}]{geometry}
\usepackage[usenames,dvipsnames,svgnames]{xcolor}
\usepackage[shortlabels]{enumitem}
\usepackage[framemethod=TikZ]{mdframed}
\usepackage{amsmath,amssymb,amsthm}
\usepackage[colorlinks]{hyperref}
\usepackage{multicol}
\usepackage{tabularx}
\usepackage{stmaryrd}

\hypersetup{
	colorlinks=true,
	linkcolor=blue,
	filecolor=magenta,
	urlcolor=magenta,
	pdftitle={MAT 534 Notes}
}

\usepackage{microtype}
\usepackage{mathtools}
\usepackage[headsepline]{scrlayer-scrpage}
\usepackage{thmtools}
\usepackage[textsize=footnotesize]{todonotes}

\mdfdefinestyle{mdbluebox}{roundcorner=10pt,innerbottommargin=9pt,
	linecolor=blue,backgroundcolor=TealBlue!5,}
\declaretheoremstyle[headfont=\sffamily\bfseries\color{MidnightBlue},
mdframed={style=mdbluebox},]{thmbluebox}

\mdfdefinestyle{mdbloobox}{frametitlefont=\bfseries,innerbottommargin=8pt,
	nobreak=true,backgroundcolor=TealBlue!5,linecolor=blue,}
\declaretheoremstyle[headfont=\bfseries\color{MidnightBlue},
mdframed={style=mdbloobox},headpunct={\\[3pt]},postheadspace=0pt,]{thmbloobox}

\mdfdefinestyle{mdredbox}{frametitlefont=\bfseries,innerbottommargin=8pt,
	nobreak=true,backgroundcolor=Salmon!5,linecolor=RawSienna,}
\declaretheoremstyle[headfont=\bfseries\color{RawSienna},
mdframed={style=mdredbox},headpunct={\\[3pt]},postheadspace=0pt,]{thmredbox}

\mdfdefinestyle{mdgreenbox}{linecolor=ForestGreen,backgroundcolor=ForestGreen!5,
	linewidth=2pt,rightline=false,leftline=true,topline=false,bottomline=false,}
\declaretheoremstyle[headfont=\bfseries\sffamily\color{ForestGreen!70!black},
mdframed={style=mdgreenbox},headpunct={ --- },]{thmgreenbox}

\mdfdefinestyle{mdorangebox}{linecolor=RedOrange,backgroundcolor=RedOrange!5,
	linewidth=2pt,rightline=false,leftline=true,topline=false,bottomline=false,}
\declaretheoremstyle[headfont=\bfseries\sffamily\color{RedOrange!70!black},
mdframed={style=mdorangebox},headpunct={ --- },]{thmorangebox}

\mdfdefinestyle{mdblackbox}{linecolor=black,backgroundcolor=RedViolet!5!gray!5,
	linewidth=3pt,rightline=false,leftline=true,topline=false,bottomline=false,}
\declaretheoremstyle[mdframed={style=mdblackbox}]{thmblackbox}

\declaretheoremstyle[spaceabove=6pt,spacebelow=6pt]{basehead}

\declaretheorem[style=thmbluebox,name=Theorem,numberwithin=section]{theorem}
\declaretheorem[style=thmredbox,name=Example,sibling=theorem]{example}
\declaretheorem[style=thmbluebox,name=Corollary,sibling=theorem]{corollary}
\declaretheorem[style=thmbluebox,name=Proposition,sibling=theorem]{prop}
\declaretheorem[style=thmbluebox,name=Lemma,numberwithin=theorem]{lemma}
\declaretheorem[style=thmbluebox,name=Theorem,numbered=no]{theorem*}
\declaretheorem[style=thmbluebox,name=Lemma,numbered=no]{lemma*}
\declaretheorem[style=thmgreenbox,name=Claim,numberwithin=theorem]{claim}
\declaretheorem[style=thmgreenbox,name=Claim,numbered=no]{claim*}
\declaretheorem[style=thmblackbox,name=Remark,numberwithin=theorem]{remark}
\declaretheorem[style=thmblackbox,name=Remark,numbered=no]{remark*}
\declaretheorem[style=thmgreenbox,name=Definition,sibling=theorem]{definition}
\declaretheorem[style=thmgreenbox,name=Definition,numbered=no]{definition*}
\declaretheorem[style=thmorangebox,name=Warning,sibling=theorem]{warning}
\declaretheorem[style=thmorangebox,name=Warning,numbered=no]{warning*}
\declaretheorem[style=thmblackbox,name=Exercise,sibling=theorem]{exercise}
\declaretheorem[style=thmblackbox,name=Exercise,numbered=no]{exercise*}

\addtolength{\textheight}{3.14cm}
\providecommand{\ol}{\overline}
\providecommand{\normaleq}{\trianglelefteq}
\providecommand{\normal}{\triangleleft}
\providecommand{\ceil}[1]{\left\lceil #1 \right\rceil}
\providecommand{\floor}[1]{\left\lfloor #1 \right\rfloor}
\newcommand{\dd}{\mathrm{d}}
\providecommand{\ul}{\underline}
\providecommand{\eps}{\varepsilon}
\providecommand{\half}{\frac{1}{2}}
\providecommand{\dang}{\measuredangle} %% Directed angle
\providecommand{\CC}{\mathbb C}
\providecommand{\FF}{\mathbb F}
\providecommand{\NN}{\mathbb N}
\providecommand{\QQ}{\mathbb Q}
\providecommand{\RR}{\mathbb R}
\providecommand{\ZZ}{\mathbb Z}
\providecommand{\dg}{^\circ}
\providecommand{\ii}{\item}
\providecommand{\setdiff}{\smallsetminus}
\providecommand{\alert}[1]{{\sffamily\textbf{\textcolor{blue}{#1}}}}
\providecommand{\inv}{^{-1}}
\providecommand{\mb}[1]{\mathbf{#1}}

\reversemarginpar

\newenvironment{soln}{\begin{proof}[Solution]}{\end{proof}}
\newenvironment{walkthrough}{\noindent \textbf{Walkthrough.}}{\medskip}
\newlist{walk}{enumerate}{3}
\setlist[walk]{label=\bfseries (\alph*)}

\providecommand{\pow}{\operatorname{Pow}}
\providecommand{\vocab}[1]{{\textbf{\textcolor{ForestGreen}{#1}}}}
\providecommand{\lcm}{\operatorname{lcm}}
\providecommand{\sgn}{\operatorname{sgn}}
\providecommand{\abs}[1]{\left|#1\right|}

\usepackage{asymptote}
\begin{asydef}
	defaultpen(fontsize(10pt));
	unitsize(1cm);
	size(8cm); // set a reasonable default
	usepackage("amsmath");
	usepackage("amssymb");
	settings.tex="pdflatex";
	settings.outformat="pdf";
	// Replacement for olympiad+cse5 which is not standard
	import geometry;
	// recalibrate fill and filldraw for conics
	void filldraw(picture pic = currentpicture, conic g, pen fillpen=defaultpen, pen drawpen=defaultpen)
	{ filldraw(pic, (path) g, fillpen, drawpen); }
	void fill(picture pic = currentpicture, conic g, pen p=defaultpen)
	{ filldraw(pic, (path) g, p); }
	// some geometry
	pair foot(pair P, pair A, pair B) { return foot(triangle(A,B,P).VC); }
	pair orthocenter(pair A, pair B, pair C) { return orthocentercenter(A,B,C); }
	pair centroid(pair A, pair B, pair C) { return (A+B+C)/3; }
	// cse5 abbreviations
	path CP(pair P, pair A) { return circle(P, abs(A-P)); }
	path CR(pair P, real r) { return circle(P, r); }
	pair IP(path p, path q) { return intersectionpoints(p,q)[0]; }
	pair OP(path p, path q) { return intersectionpoints(p,q)[1]; }
	path Line(pair A, pair B, real a=0.6, real b=a) { return (a*(A-B)+A)--(b*(B-A)+B); }
	// cse5 more useful functions
	picture CC() {
		picture p=rotate(0)*currentpicture;
		currentpicture.erase();
		return p;
	}
	pair MP(Label s, pair A, pair B = plain.S, pen p = defaultpen) {
		Label L = s;
		L.s = "$"+s.s+"$";
		label(L, A, B, p);
		return A;
	}
	pair Drawing(Label s = "", pair A, pair B = plain.S, pen p = defaultpen) {
		dot(MP(s, A, B, p), p);
		return A;
	}
	path Drawing(path g, pen p = defaultpen, arrowbar ar = None) {
		draw(g, p, ar);
		return g;
	}
\end{asydef}

\usepackage{mathrsfs}

\ihead{\footnotesize MAT 534 Notes}
\ohead{\footnotesize \today}

\title{\vspace{-2em}MAT 534: Algebra I}
\author{\large Aditya Pahuja}
\date{}
\begin{document}
\maketitle
\tableofcontents

\pagebreak

\section{August 25}\label{sec:aug25}

\begin{definition}[Group]
    \label{def:group}
    A \vocab{group} is a set $G$ with a binary operation
    \begin{align*}
        *\colon G\times G &\to G \\
	(g,h)&\mapsto g*h
    \end{align*}
    such that
    \begin{enumerate}[(1)]
        \ii $*$ is associative, so $f*(g*h) = (f*g)*h$.
	\ii There is a \vocab{neutral element} $e\in G$ such that
	for all $g\in G$, $e*g=g*e=g$.
	\ii For all $g\in G$, there is an \vocab{inverse element}
	$h\in G$ such that $g*h=h*g=e$.
    \end{enumerate}
\end{definition}

Some basic facts:
\begin{itemize}
    \ii $e$ is unique, because if $e$ and $e'$ are neutral
    then $e = e*e' = e'$.
    \ii For each $g\in G$, the inverse of $g$ is unique, for
    if $h$ and $h'$ are inverses of $g$ then
    $h = h*(g*h') = (h*g)*h' = h'$.
    In a generic group we write $g\inv$ for the inverse of $g$.
    \ii The inverse of $g\inv$ is $g$.
\end{itemize}

From here on, when working with a generic group,
we use multiplicative notation, meaning
\begin{itemize}
    \ii multiplication is written either with juxtaposition or,
    if emphasis is need, a multiplication dot: $g*h = gh = g\cdot h$.
    \ii repeated multiplication is given by exponents:
    if $n$ is a positive integer, then $g^n$ is the $n$-fold product of $g$;
    if $n = 0$, then $g^0 = 1$, and if $n$ is negative, then
    $g^n = (g\inv)^{-n}$, the $-n$-fold product of $g\inv$.
\end{itemize}

\begin{definition}[Abelian group]
    \label{def:abeliangroup}
    $G$ is \vocab{abelian} if its multiplication operation is commutative,
    i.e. $gh = hg$ for all $g,h\in G$.
\end{definition}

For generic abelian groups we instead prefer additive notation, so
\begin{itemize}
    \ii we write $g+h$ instead of $gh$ or $g\cdot h$ or $g*h$ for the operation;
    \ii we write $0$ instead of $1$ for the neutral element;
    \ii we write $-g$ instead of $g\inv$ for inverses;
    \ii we write $ng$ instead of $g^n$ for $n$-fold applications of the operation
    where $n$ is an integer.
\end{itemize}

\begin{example}
    [Examples of groups]
    Some abelian groups are the additive groups
    $(\ZZ,+)$, $(\QQ,+)$, $(\RR,+)$, $(\CC,+)$, \dots
    and the multiplicative groups
    $(\QQ^\times, \cdot)$, $(\RR^\times,\cdot)$, etc.
    where the superscript $\times$ means ``nonzero''
    (so $\QQ^\times$ is the set of nonzero rationals.)
\end{example}

In real life groups often show up as symmetries of objects
(where ``symmetry'' is often interpreted as ``structure-preserving,''
such as, say, the group of isomorphisms between two groups,
the group of automorphisms on an algebraic extension of $\QQ$,
the group of deck transformations on a covering space, etc.)
A fairly literal example of this is the \vocab{dihedral group} $D_{2n}$,
which is the group of symmetries of a regular $n$-gon, i.e.
the set of rigid motions in space that map an $n$-gon onto itself,
like rotations by $\frac{2\pi}n$ about the center
or reflections over certain lines; the group operation
is performing these rigid motions one after the other.

\begin{exercise}
    Convince yourself that these transformations form a group.
\end{exercise}

To abstract this into a more symbolic form,
suppose that the vertices of the $n$-gon are the points $e^{2\pi k/n}$
for $k = 1,\dots,n$ in the complex plane, so that the real line
is an axis of symmetry of the $n$-gon. Then, let $r$ represent
rotation of the $n$-gon counterclockwise by $\frac{2\pi}n$
and $d$ reflection over the real axis.
\begin{exercise}
    Show that $drd = r\inv$, and hence that the elements of $D_{2n}$ are
    the powers of $r$ ($1$, $r$, \dots, $r^{n-1}$)
    and $d$ times these powers of $r$ ($d$, $dr$, \dots, $dr^{n-1}$).
    This means that $D_{2n}$ has $2n$ elements, justifying the subscript.
\end{exercise}

We can also express $D_{2n}$ with a \vocab{group presentation}, namely
\begin{equation*}
    D_{2n} = \langle r,d \mid r^n = 1,\, d^2 = 1,\, drd = r\inv\rangle;
\end{equation*}
this says that $D_{2n}$ is the group with two generators $r$, $d$
that can be algebraically manipulated with the given relations.
This is in general not that fruitful of a task, though,
since the problem of deciding whether two elements of a group are equal
by looking at the group's presentation is undecidable.

Another example is the \vocab{matrix group} $GL_n(F)$,
which is the group of invertible $n\times n$ matrices
whose entries are in the field $F$ (this of course consists of
the $n\times n$ matrices $A$ with nonzero determinant),
whose operation is matrix multiplication. Checking that this is a group:
\begin{enumerate}[(1)]
    \ii Matrix multiplication is associative because
    matrices in $GL_n(F)$ are just concrete representations
    of linear maps on $F^n$, and function composition is associative.
    \ii The neutral element is the identity matrix
    \begin{equation*}
        I = \begin{pmatrix}
	    1&0&\cdots&0\\
	    0&1&\cdots&0\\
	    \vdots&\vdots&\ddots&\vdots\\
	    0&0&\cdots&1\\
        \end{pmatrix}.
    \end{equation*}
    \ii The inverse of a matrix can either be computed concretely
    or, again understanding matrices to be linear maps,
    one can directly construct the inverse of the corresponding linear map
    by looking at bases of the vector spaces that the map acts on.
\end{enumerate}

A final very useful group is the \vocab{symmetric group} $S_n$,
which is the group of \vocab{permutations} on a set of $n$ things
(for concreteness, the $n$ things are the elements of $[n] := \{1,\dots,n\}$);
that is, $S_n$ is the set of bijections $\sigma\colon[n]\to[n]$
endowed with the operation of function composition
(thus multiplication happens right-to-left here;
we did not need to worry about this with matrices because
we can imagine them as just grids of numbers
for which we have an admittedly obtuse multiplication rule).
Note that $|S_n| = n!$, since there are $n$ choices for $\sigma(1)$,
followed by $n-1$ choices for $\sigma(2)$, and so on,
for a total of $n\cdot(n-1)\cdots2\cdot 1 = n!$.

When working with permutations, we would like to have a good way to write them down;
naively, one can simply write down a table
\begin{equation*}
    \begin{tabular}{c|c|c|c|c|c|c|c|c|c|c|c|c|c}
	$i$ & 1 & 2 & 3 & 4 & 5 & 6 & 7 & 8 & 9 & 10 & 11 & 12 & 13 \\
	\hline
	$\sigma(i)$ & 12 & 13 & 3 & 1 & 11 & 9 & 5 & 10 & 6 & 4 & 7 & 8 & 2
    \end{tabular}
\end{equation*}
but this is clunky and annoying, as well as
not telling us much about the structure of $\sigma$
(which is why we'd rather not just resort to a more compact 13-tuple).

Instead, we use \vocab{cycle notation}.
\begin{definition}[Cycles]
    \label{def:Sn-cycles}
    Given distinct $a_1,\dots,a_k\in\{1,\dots,n\}$, define
    \begin{equation*}
	(a_1\; a_2\; \dots\; a_k)
    \end{equation*}
    to be the permutation that maps
    $a_1\mapsto a_2\mapsto\cdots\mapsto a_k\mapsto a_1$.
    We call this a \vocab{$k$-cycle}.
\end{definition}

It is a theorem (that we will not prove rigorously) that
any permutation in $S_n$ can uniquely by expressed as a product
$\sigma_1\sigma_2\cdots\sigma_r$ of pairwise disjoint nontrivial cycles $\sigma_j$,
up to the ordering of the $\sigma_j$
(pairwise disjoint every element of $[n]$ is fixed by
all but one of the $\sigma_j$, and nontrivial meaning
the cycles aren't the identity permutation, so not $1$-cycles).
The way that one computes this cycle representation of a permutation
(and, incidentally, proves that this representation exists and is unique)
is simply by following where the numbers are mapped to, say by using the table.
For example, in the table of $\sigma\in S_{13}$ above,
\begin{itemize}
    \ii $1\mapsto 12\mapsto 8\mapsto 10\mapsto 4\mapsto 1$,
    giving the cycle $(1\; 12\; 8\; 10\; 4)$.
    \ii $2\mapsto13\mapsto 2$,
    giving the cycle $(2\;13)$.
    \ii $3$ maps to itself,
    giving the cycle $(3)$, which we omit.
    \ii $5\mapsto 11\mapsto 7$,
    giving the cycle $(5\;11\;7)$.
    \ii $6\mapsto 9\mapsto 6$,
    giving the cycle $(6\;9)$.
\end{itemize}
Therefore,
\begin{equation*}
    \sigma = (1\; 12\; 8\; 10\; 4)(2\;13)(5\;11\;7)(6\;9).
\end{equation*}

As another example of cycle notation, we can write
\begin{equation*}
    S_3 = \{(1), (1\; 2), (1\; 3), (2\; 3), (1\; 2\; 3), (1\; 3\; 2)\}
\end{equation*}
where by convention the identity permutation is written as $(1)$
because the empty string is hard to write.
Recalling that composition of permutations happens right-to-left,
we can compute
\begin{equation*}
    (1\; 2)(1\; 3) = (1\; 3\; 2)
\end{equation*}
by following where each element goes (imagine moving the numbers around in a line):
$1$ goes to position $3$, and then stays there because $(1\; 2)$
doesn't touch position $3$; the newly displaced $3$ goes to position $1$
and then is promptly moved to position $2$; and finally $2$,
which is not moved by $(1\;3)$, is sent to position $1$,
so the overall map is $1\mapsto3\mapsto 2\mapsto1$.
On the other hand, $(1\;3)(1\;2) = (1\; 2\; 3)$ by the same logic,
so in fact we have shown that $S_3$
(and generally $S_n$)
is non-abelian for all $n\ge 3$.

\section{August 27}\label{sec:aug27}

\begin{definition}[Homomorphism]
    \label{def:homomorphism}
    Let $G$ and $H$ be groups. A \vocab{homomorphism} $\varphi\colon G\to H$
    is a function such that
    \begin{equation*}
        \varphi(g_1g_2) = \varphi(g_1)\varphi(g_2).
    \end{equation*}
    If $\varphi$ is bijective, it is called an \vocab{isomorphism}
    and $G$, $H$ are said to be \vocab{isomorphic},
    which is denoted as $G\cong H$.
\end{definition}

Some simple facts:
\begin{itemize}
    \ii $\varphi(1_G) = 1_H$ because $\varphi(1_G) = \varphi(1_G)\cdot\varphi(1_G)$.
    \ii $\varphi(g\inv) = \varphi(g)\inv$ because
    $\varphi(g)\varphi(g\inv) = \varphi(1_G) = 1_H$.
    \ii If $\varphi\colon G\to H$ is an isomorphism
    then so is $\varphi\inv\colon H\to G$.
\end{itemize}

In fact, isomorphism is an equivalence relation on groups.

\begin{example}
    Some examples of homomorphisms:
    \begin{enumerate}[(1)]
	\ii The determinant $\det\colon GL_n(F)\to F^\times$
	is multiplicative, so it defines a homomorphism
	on the set of invertible linear maps.
	\ii If the vertices of a regular $n$-gon are labeled
	$1$, \dots, $n$, then $D_{2n}$ permutes these vertices;
	therefore we get a homomorphism $D_{2n}\to S_n$
	by sending $d\mapsto (1)(2\; n)(3\; n-1)\cdots$
	and $r\mapsto (1\; 2\; \cdots \; n)$.
	\ii The groups $(\RR,+)$ and $(\RR_{>0}, \cdot)$
	are isomorphic, with isomorphism $\exp\colon\RR\to\RR_{>0}$.
	\ii Fix $g\in G$. Then, the \vocab{conjugation map} is an isomorphism
	$\varphi_g\colon G\to G$ given by $\varphi_g(x) = gxg\inv$.
	Moreover, $\varphi_a\varphi_b = \varphi_ab$,
	so the group $\{\varphi_g : g\in G\}$ with operation being group composition
	is isomorphic to $G$.
    \end{enumerate}
\end{example}

\begin{definition}[Subgroup]
    \label{def:subgroup}
    A \vocab{subgroup} of a group $G$ is a subset $H\subseteq G$
    that is a group under the operation of $G$; in particular,
    \begin{enumerate}[(1)]
        \ii $1_G\in H$
	\ii $h_1,h_2\in H\implies h_1h_2\in H$
	\ii $h\in H\implies h\inv\in H$.
    \end{enumerate}
    We more succinctly write $H\le G$ if $H$ is a subgroup of $G$.
    Furthermore $H$ is \vocab{normal} in $G$
    if $\varphi_g(H) = H$, i.e. $h\in H$ implies $ghg\inv\in H$;
    this is denoted $H\trianglelefteq G$.
\end{definition}

Normal subgroups are a very important part
of the structure of a group $G$ due to the following proposition:

\begin{prop}
    \label{prop:kernel-is-normal}
    If $\varphi\colon G\to H$ is a homomorphism
    then $\varphi(G)$ is a subgroup of $H$
    and the kernel $\ker\varphi = \{g\in G : \varphi(G) = 1_H\}$
    is a normal subgroup of $G$.
\end{prop}

\begin{proof}
    Trivial; omitted.
\end{proof}

\begin{example}
    Some examples:
    \begin{enumerate}[(1)]
        \ii $SL_n(F)$, the subset of $GL_n(F)$
	whose determinant is $1$,
	is the kernel of the determinant homomorphism,
	so it is normal in $GL_n(F)$.
	\ii The \vocab{center} $Z(G) = \{g\in G : gh = hg\forall h\in G\}$
	is a normal subgroup of $G$.
    \end{enumerate}
\end{example}

\begin{prop}
    \label{prop:subgrp-intersection}
    The intersection of an arbitrary collection of subgroups
    of a group $G$ is also a subgroup of $G$.
\end{prop}

\begin{proof}
    Trivial; omitted.
\end{proof}

\begin{definition}[Generated subgroup]
    \label{def:generated-subgrp}
    Let $A$ be a subset of a group $G$.
    The subgroup $\langle A\rangle$ generated by $A$
    is the smallest (with respect to inclusion)
    subgroup of $G$ containing $A$, or equivalently
    the set of elements of $G$ that can be expressed as
    a finite product of elements in $\{a,a\inv : a\in A\}$
    (i.e. take $\{a,a\inv : a\in A\}$ and shake vigorously).
\end{definition}

To be precise,
\begin{equation*}
    \langle A\rangle = \bigcap_{\substack{H\supseteq A\\ H\le G}} H.
\end{equation*}
This is not too practical for concrete descriptions of generated subgroups,
but still good to keep in mind.

\begin{example}
    Given $g\in G$, the generated subgroup $\langle g\rangle$
    is $\{g^n : n\in\ZZ\}$. If this set is infinite,
    then it is isomorphic to $\ZZ$ via the isomorphism $n\mapsto g^n$.
    In this case, we say that $g$ has infinite order.
    Otherwise, $g^m = 1$ for some $m$; if $m$ is taken to be
    positive and as small as possible, then
    the kernel of the $n\mapsto g^n$ isomorphism has kernel
    $m\ZZ = \{mn : n\in\ZZ\}$, hence $\langle g\rangle \cong \ZZ/m\ZZ$
    (which is the group of integers modulo $m$ under addition).
    We say here that $g$ has order $m$.
    In general we denote the order of $g\in G$ with $\abs g$.
\end{example}

\begin{theorem}[Lagrange]
    \label{thm:lagrange-thm}
    If $G$ is finite and $H$ is a subgroup of $G$,
    then $\abs H$ divides $\abs G$.
\end{theorem}

\begin{proof}
    For $a,b\in G$, say that $a\sim b$ if $a\inv b \in H$.
    It is easy to see that this an equivalence relation
    due to $H$ being a subgroup of $G$. For a fixed $a$, then,
    the equivalence class of $a$
    \begin{equation*}
	\{b\in G : b\sim a\} = \{ah : h\in H\} \eqcolon aH
    \end{equation*}
    (also called a \vocab{left coset} of $H$)
    has the same size as $H$, hence
    $\abs G = \abs H\cdot(\text{\# of left cosets of $H$})$
    so $\abs G/\abs H$ is an integer.
\end{proof}

The number of left cosets of $H$ in $G$ is called the \vocab{index}
of $H$ in $G$ and is denoted $(G:H)$; the proof above shows that
if $G$ is finite then $(G:H) = \abs G/\abs H$.

Another important homomorphism is the \vocab{sign homomorphism}
$\sgn\colon S_n\to\{\pm1\}$. Noting that $S_n$ is generated by
the set of $2$-cycles (\vocab{transpositions}),
we say that $\sgn((a\; b)) = -1$ for all $a\ne b$ in $[n]$
and then extend multiplicatively
(so for example $3$-cycles, which are always a product of two $2$-cycles,
all return $1$ when passed through the $\sgn$ homomorphism).
The permutations $\sigma$ with $\sgn(\sigma) = -1$
are called odd permutations and those with $\sgn(\sigma) = 1$ are even.
The even permutations, being the kernel of $\sgn$,
form a normal subgroup of $S_n$ called the \vocab{alternating group} $A_n$.

\section{September 3}\label{sec:sep03}

If $N\trianglelefteq G$ then
\begin{equation*}
    Nb = bN \implies aN\cdot bN = abN,
\end{equation*}
i.e. multiplication of cosets is well-defined here.
\begin{definition}[Quotient group]
    \label{def:quotient-group}
    Given $N\normaleq G$, the \vocab{quotient group} or \vocab{factor group}
    $G/N$ is the group of left cosets of $N$
    with the multiplication $aN\cdot bN = abN$.
\end{definition}

Note that Lagrange's theorem tells us $\abs{G/N} = (G:N)$,
the index of $N$ in $G$.
Additionally, there is a natural homomorphism $\pi\colon G\to G/N$ mapping
\begin{equation*}
    \pi(g) = gN,
\end{equation*}
whose kernel is $\{g\in G : gN = N\} = N$, so every normal subgroup of $G$
is the kernel of some homomorphism out of $G$.

Now we turn to the four isomorphism theorems,
which tell us about the structure of a group's subgroups
in terms of homomorphisms going out of the group.

\begin{theorem}[First isomorphism theorem]
    \label{thm:first-iso-thm}
    Let $\varphi\colon G\to H$ be a group homomorphism with kernel $K$.
    Then, $K\normaleq G$ and $\varphi(G) = \ol G$ is
    a subgroup of $H$ isomorphic to $G/K$ with isomorphism
    $G/K\to\ol G$ given by $gK\mapsto \varphi(g)$.
\end{theorem}

This is a very simple and powerful theorem;
proving that a group is isomorphic to a quotient
often happens via the first isomorphism theorem.

\begin{theorem}[Second isomorphism theorem]
    \label{thm:second-iso-thm}
    Let $G$ be a group and $H\le G$, $K\normaleq G$.
    \begin{enumerate}[(1)]
	\ii $HK = \{h\cdot k : h\in H, k\in K\}$ is a subgroup of $G$.
	\ii $H\cap K\normaleq H$.
	\ii $HK/K\cong H/(H\cap K)$.
    \end{enumerate}
\end{theorem}

\begin{example}
    Consider the normal subgroup $K = \{(1), (1\;2)(3\;4), (1\;3)(2\;4), (1\;4)(2\;3)\}$
    of $S_4$ and $H$ the subgroup of permutations that fix $4$, so $H\cong S_3$.
    Then, $H\cap K = \{1\}$, so by \autoref{thm:second-iso-thm}
    \begin{equation*}
        HK/K \cong H/H\cap K = H \cong S_3.
    \end{equation*}
    This implies that $\abs{HK} = \abs K\cdot \abs{S_3} = 24$,
    so in fact $HK = S_4$.
\end{example}

The third and fourth isomorphism theorems are very closely related:
\begin{theorem}[Third isomorphism theorem]
    \label{thm:third-iso-thm}
    Let $G$ be a group and $N\normaleq G$. If $K\normaleq G$ is
    such that $N\le K\le G$ then
    \begin{enumerate}[(1)]
        \ii $K/N\normaleq G/N$
	\ii $(G/N)/(K/N)\cong G/K$.
    \end{enumerate}
\end{theorem}

This one says that normal subgroups $K$ ``pass through'' quotients
as long as $K$ contains the group $N$ that $G$ is being quotiented by.

\begin{theorem}[Fourth isomorphism theorem]
    \label{thm:fourth-iso-thm}
    Suppose $G$ is a group and $N\normaleq G$.
    There is a 1-1 correspondence between subgroups $H\le G$ containing $N$
    and subgroups of $G/N$ by the mapping $H\mapsto H/N$.
\end{theorem}

The fourth isomorphism theorem says that in fact the subgroups
which pass through the quotient $G/N$
(in the same spirit as in \autoref{thm:third-iso-thm})
are all the subgroups of $G/N$.
If one were to look at the subgroups of $G$ as a poset
ordered by inclusion, the fourth isomorphism theorem says that
the corresponding poset for $G/N$ looks like zooming in on
just the section of the poset bigger than $N$ and smaller than $G$.

Thus, if $N\normaleq G$, then we can try to understand
the structure of $G$ by examining
\begin{enumerate}[(1)]
    \ii the smaller groups $N$ and $G/N$.
    Repeating this analysis eventually reduces to the case where
    $G$ has no nontrivial normal subgroups, where $G$ is called \vocab{simple}.
    \ii how $G$ is put together from $N$ and $G/N$.
\end{enumerate}

\section{September 8}\label{sec:sep08}

Two ways to study a group:
\begin{enumerate}[(1)]
    \ii For a field $F$ consider $GL_n(F)$.
    We can study $G$ by homomorphisms $\rho\colon G\to GL_n(F)$;
    this is called \vocab{representation theory}
    (not discussed in this class)..
    \ii For a set $X$, we can consider the symmetric group
    $S_X = \{f\colon X\to X : f\text{ is bijective}\}$.
    Then, we can study $G$ by homomorphisms $\rho\colon G\to S_X$.
    In this case we say that $G$ \vocab{acts on} $X$
    or that $G$ defines a \vocab{group action} on $X$.
\end{enumerate}

We shall elaborate on the latter point.
\begin{definition}[Group action]
    \label{def:group-action}
    For each $g\in G$, we associate a bijection $\rho(g)\colon X\to X$
    and write $g\cdot x \coloneq \rho(g)(x)$ such that the following hold:
    \begin{enumerate}[(1)]
	\ii $\rho(gh) = \rho(g)\circ \rho(h)$
	\ii $\rho(1) = \operatorname{id}_X$, so $1\cdot x = x$ for all $x$.
    \end{enumerate}
    As mentioned above, we say that $G$ \vocab{acts on} the set $X$.
    Since the group action is written to the left here,
    this is called a \vocab{left action} of $G$ on $X$.
\end{definition}

\begin{example}
    Some examples of group actions:
    \begin{enumerate}[(1)]
	\ii $D_8$ acts on the corners of a square,
	as per its definition. If one were to label the vertices
	$1$ through $4$, we would get a homomorphism $D_8\to S_4$
	(see example [REFERENCE, 2.2]).
	\ii If $G$ is a group and $H$ a subgroup,
	then $G$ acts on the set of all left cosets of $H$
	by $g\cdot aH = gaH$; it is an action because
	multiplication in $G$ is associative.
	If $(G:H)$ is finite then this gives a homomorphism
	$G\to S_n$.
	\ii A group $G$ can act on itself via conjugation, namely
	$g\cdot h = ghg\inv$.
    \end{enumerate}
\end{example}

\begin{definition}[Orbits and stabilizers]
    \label{def:orbit-stab}
    Suppose $G$ acts on a set $X$. For a given $x\in X$,
    the \vocab{orbit} of $x$ is the set $G\cdot x = \{g\cdot x : g\in G\}$.
    An element $x\in X$ is a \vocab{fixed point} if it is fixed by all $g\in G$,
    i.e. $g\cdot x = x$ for all $g$.
    $G$ acts \vocab{transitively} on $X$ if all of $X$ is one orbit.
    Given $x\in X$, the \vocab{stabilizer} of $x$ is
    the subgroup of $G$ that fixes $x$, namely
    \begin{equation*}
	G_x \coloneq \{g\in G : g\cdot x = x\}.
    \end{equation*}
\end{definition}

\begin{exercise}
    Check that ``being in the same orbit'' is an equivalence relation
    on $X$, hence the orbits (equivalence classes) partition $X$.
\end{exercise}

\begin{exercise}
    Show that for any $x\in X$, $G_x\le G$.
\end{exercise}

\begin{example}
    [A famous example]
    Let $\mathbb H = \{z\in\CC : \operatorname{Im}(z)>0\}$
    be the upper half plane and $G = SL_2(\RR)$.
    Then, $G$ acts on the Riemann sphere $\hat C = \CC\cup\{\infty\}$ by
    \begin{equation*}
        \begin{pmatrix}
	    a&b\\c&d
        \end{pmatrix} \cdot z =
	\frac{az+b}{cz+d}.
    \end{equation*}
    Verifying that this is an action is a simple calculation.
    
    The orbit of $0$ is the entire real line $\RR\cup\{\infty\}$.
    % TODO
\end{example}

If $y\in G\cdot x$, say $y = g\cdot x$, then
\begin{equation*}
    h\in G_y\implies h\cdot y = y \implies
    h\cdot g\cdot x = g\cdot x \implies (g\inv hg)\cdot x = x,
\end{equation*}
so if $x$ is in the orbit of $y$ then elements of the stabilizer of $y$
are conjugate to elements of the stabilizer of $x$;
specifically, $G_y = g\cdot G_x \cdot g\inv$.

Consider $f\colon G\to G\cdot x$ by $g\mapsto g\cdot x$.
Obviously $f$ is surjective. Moreover, consider
\begin{equation*}
    f\inv(y) = \{g\in G : g\cdot x = y\}.
\end{equation*}
If $g\cdot x = h\cdot x$, then $h\inv g\cdot x = x$, meaning
$h\inv g\in G_x$ or equivalently $g$ is in the left coset $h\cdot G_x$.
This means $f\inv(y)$ is a left coset of $G_x$,
hence left cosets of $G_x$ are in bijection with the orbit $G\cdot x$
via $gG_x\mapsto g\cdot x$. In the case where the orbit is finite,
the size of the orbit is $\abs{G\cdot x} = (G : G_x)$.
When $G$ is a finite group acting on a set $X$, we thus have
\begin{equation*}
    \abs G = \abs{G\cdot x}\cdot \abs{G_x}.
\end{equation*}
When $X$ is also finite, we have $X$ as the disjoint union $\bigsqcup G\cdot x_i$
where we choose one representative $x_i$ of each orbit, so
\begin{equation*}
    \abs X = \sum_{i} \abs{G\cdot x_i} = \sum_i \frac{\abs G}{\abs{G_{x_i}}}.
\end{equation*}

An important special case is the conjugation action described in
[REFERENCE whichever example i wrote about conjugation in],
namely $\alpha\colon G\times G\to G$ by $\alpha(g,h) = ghg\inv$.
The orbit of an $h\in G$ is the \vocab{conjugacy class}
$\{ghg\inv : g\in G\}$ of $h$.
The stabilizer of $h$ is
\begin{equation*}
    G_h = \{g\in G : ghg\inv = h \iff gh = hg\},
\end{equation*}
i.e. the elements in $G$ that commute with $h$.
This set is called the \vocab{centralizer} $C(h)$ of $h$.
Of course, this means we can write the center of $G$ as
\begin{equation*}
    Z(G) = \bigcap_{h\in G} C(h).
\end{equation*}

\begin{exercise}
    When is $h\in G$ a fixed point under conjugation?
\end{exercise}

If $G$ is finite and written as the union of conjugacy classes of
$h_i$ (so we pick one $h_i$ from each conjugacy class)
\begin{equation*}
    \abs G = \sum_i \abs{\text{conjugacy class of $h_i$}}
    = \sum_i \frac{\abs G}{\abs{C(h_i)}}
    = \abs{Z(G)} + \sum_{C(h_j)\ne G} \frac{\abs G}{\abs{C(h_j)}}
\end{equation*}
where the last expression groups together all the singleton orbits
and so the summation runs over representatives $h_j$
of orbits with size greater than $1$ (i.e. $h_j\notin Z(G)$).
This is called the \vocab{class equation} of $G$.
The neat thing is that analyzing the addends in this equation
can give constraints on the group
because each adddend must be a factor of $\abs G$.










\end{document}

